\documentclass[11pt]{article}
\usepackage{coursemacros}

\title{Probability --- Exam 1 Practice\\[2pt]\large Chapters 1--3}
\author{Mikey Ferguson}
\date{}

\begin{document}
\maketitle

\noindent\textbf{Part I is a full exam: 75 minutes, open notes, no calculator.}
Leave arithmetic unsimplified wherever it is right: $\tfrac{1}{3}\cdot\tfrac{3}{10}$,
$(5/6)^{10}$, $e^{-2}$ and $\Phi(1)$ are all finished answers. Part II is extra
practice on what Part I does not reach; it is not timed.

\section*{Part I --- timed, 75 minutes}

\begin{problem}{1}
  (Conditioning and Bayes.) Box 1 holds 3 red and 2 blue balls. Box 2 holds 2 red
  and 3 blue balls. Roll a fair die: on a 1 or 2 use Box 1, otherwise Box 2. Draw
  two balls from that box without replacement.
  \begin{enumerate}
    \item[(a)] Find the probability that both balls are red.
    \item[(b)] Given that both balls are red, find the probability that Box 1 was used.
    \item[(c)] Using only the axioms of probability, prove that if $E \subseteq F$ then
      $P(E) \le P(F)$.
  \end{enumerate}
\end{problem}
\vfill\newpage

\begin{problem}{2}
  (Indicators.) A fair die is rolled 10 times.
  \begin{enumerate}
    \item[(a)] Let $X$ be the number of different faces that appear at least once.
      Find $\EE{X}$.
    \item[(b)] Let $Y$ be the number of faces that appear exactly once. Find $\EE{Y}$.
  \end{enumerate}
\end{problem}
\vfill\newpage

\begin{problem}{3}
  (A joint density.) $X$ and $Y$ have joint density
  \[
    f(x,y) = 2, \qquad 0 < y < x < 1,
  \]
  and $0$ elsewhere.
  \begin{enumerate}
    \item[(a)] Find the marginal densities $f_X$ and $f_Y$.
    \item[(b)] Find $\Cov(X,Y)$.
    \item[(c)] Find $P(X + Y < 1)$.
    \item[(d)] Find the conditional density of $Y$ given $X = x$, and $\EEc{Y}{X}$.
    \item[(e)] Use the conditional variance formula
      $\Var(Y) = \EE{\Var(Y \mid X)} + \Var\!\left(\EEc{Y}{X}\right)$ to find
      $\Var(Y)$. Check it against $f_Y$.
  \end{enumerate}
\end{problem}
\vfill\newpage

\begin{problem}{4}
  (Moment generating functions.)
  \begin{enumerate}
    \item[(a)] Show that a $\Poisson{\lambda}$ random variable has moment generating
      function $\mgf{}(t) = e^{\lambda(e^t - 1)}$.
    \item[(b)] A random variable $X$ has moment generating function
      \[
        \mgf{X}(t) = \left(\tfrac{1}{3} + \tfrac{2}{3}e^{t}\right)^{5} e^{2(e^{t} - 1)}.
      \]
      Write $X$ as a sum of two independent random variables with named
      distributions, and justify the step that lets you do so.
    \item[(c)] Find $\EE{X}$ and $\Var(X)$.
    \item[(d)] Find $P(X = 0)$ and $P(X = 1)$.
  \end{enumerate}
\end{problem}
\vfill\newpage

\begin{problem}{5}
  (Limit theorems.) A lamp burns one bulb at a time and a dead bulb is replaced
  at once. Bulb lifetimes are independent with mean 10 hours and standard
  deviation 10 hours. Let $S$ be the total burning time of 36 bulbs.
  \begin{enumerate}
    \item[(a)] Use Markov's inequality to bound $P(S \ge 720)$.
    \item[(b)] Use Chebyshev's inequality to bound $P(|S - 360| \ge 120)$.
    \item[(c)] Use the central limit theorem to approximate $P(S > 420)$, in terms
      of $\Phi$.
  \end{enumerate}
\end{problem}
\vfill\newpage

\begin{problem}{6}
  (Random sums.) The number of customers entering a shop in an hour is
  $N \dist \Poisson{4}$.
  \begin{enumerate}
    \item[(a)] Each customer buys something with probability $\tfrac14$,
      independently of everything else. Let $B$ be the number who buy. Show, by
      conditioning on $N$, that $B \dist \Poisson{1}$.
    \item[(b)] Customer $i$ spends $Y_i$ dollars, where $Y_1, Y_2, \ldots$ are iid,
      independent of $N$, with mean 3 and variance 2. Let
      $T = \sum_{i=1}^{N} Y_i$. Find $\EE{T}$ and $\Var(T)$.
  \end{enumerate}
\end{problem}
\vfill\newpage

\section*{Part II --- extra practice}

\begin{problem}{7}
  (Which comes first.) Two fair dice are rolled repeatedly. Find the probability
  that a sum of 5 appears before a sum of 7.
\end{problem}
\vspace{4cm}

\begin{problem}{8}
  (A joint table.) $X$ and $Y$ have joint mass function
  \[
    \begin{array}{c|ccc}
          & Y = 0 & Y = 1 & Y = 2 \\ \hline
    X = 0 & 1/8   & 1/4   & 1/8   \\
    X = 1 & 1/4   & 1/8   & 1/8
    \end{array}
  \]
  \begin{enumerate}
    \item[(a)] Are $X$ and $Y$ independent?
    \item[(b)] Find $\EEc{X}{Y = y}$ for each $y$, and check that
      $\EE{\EEc{X}{Y}} = \EE{X}$.
    \item[(c)] Find $\Cov(X,Y)$.
  \end{enumerate}
\end{problem}
\vfill\newpage

\begin{problem}{9}
  (Change of variables.) $X$ and $Y$ are independent $\Expo{1}$. Let
  $U = X + Y$ and $V = X/(X+Y)$. Find the joint density of $U$ and $V$. Are $U$
  and $V$ independent? Name their distributions.
\end{problem}
\vspace{5cm}

\begin{problem}{10}
  (Hypergeometric by indicators.) An urn holds 6 white and 4 black balls. Draw 5
  without replacement and let $S$ be the number of white balls drawn. Write $S$
  as a sum of indicators and find $\EE{S}$ and $\Var(S)$ from them.
\end{problem}
\vfill\newpage

\begin{problem}{11}
  (First-step analysis.) A fair die is rolled until the first 6. Let $S$ be the
  sum of all the rolls, the 6 included. Find $\EE{S}$ by conditioning on the first
  roll.
\end{problem}
\vspace{4cm}

\begin{problem}{12}
  (Splitting a Poisson sum.) $X \dist \Poisson{\lambda}$ and $Y \dist \Poisson{\mu}$
  are independent.
  \begin{enumerate}
    \item[(a)] Use moment generating functions to show $X + Y \dist \Poisson{\lambda + \mu}$.
    \item[(b)] Find the conditional distribution of $X$ given $X + Y = n$.
  \end{enumerate}
\end{problem}
\vspace{4cm}

\begin{problem}{13}
  (No mean.) $X$ has the Cauchy density $f(x) = \dfrac{1}{\pi(1+x^2)}$. Show that
  $\EE{|X|} = \infty$, so $\EE{X}$ does not exist.
\end{problem}
\vfill

\end{document}
